\subsection{Execution library}

\begin{frame}[t,fragile]{Senders and receivers}
  \begin{itemize}
    \item The execution library provides a set of concepts and utilities for asynchronous programming.

    \mode<presentation>{\vfill\pause}
    \item \textemph{Key concepts}:
      \begin{itemize}
        \item \textgood{Senders}: Objects that represent asynchronous operations.
        \item \textgood{Receivers}: Objects that receive completion signals from an asynchronous operation.
        \item \textgood{Schedulers}: Objects that represent execution contexts for asynchronous operations.
      \end{itemize}

    \mode<presentation>{\vfill\pause}
    \item \textemph{Advantages}:
      \begin{itemize}
        \item \textmark{Structured concurrency}: Clear begin and end points for asynchronous operations.
        \item \textmark{Decoupling} the \textgood{computation} from the \textgood{execution context}.
        \item \textmark{Decoupling} the \textgood{initiation} of asynchronous operations from their \textgood{completion} handling.
        \item \textmark{High performance}: Can avoid heap allocations, depending on the sender chain and implementation.
        \item \textmark{Universal composition}: Flexible composition of computation.
      \end{itemize}
  \end{itemize}
\end{frame}

\begin{frame}[t,fragile]{Execution framework}
  \begin{itemize}
    \item \textmark{Schedulers}
      \begin{itemize}
        \item Handle to a computation resource (e.g., thread pool, event loop).
        \item Can be used to specify where an asynchronous operation should execute.
      \end{itemize}

    \mode<presentation>{\vfill\pause}
    \item \textmark{Senders}
      \begin{itemize}
        \item Represent asynchronous operations.
        \item Can be chained with other senders to create complex asynchronous workflows.
      \end{itemize}

    \mode<presentation>{\vfill\pause}
    \item \textmark{Receivers}
      \begin{itemize}
        \item Handle the completion of asynchronous operations.
        \item Can be invoked with a result, an error, or a stopped/cancelled signal.
      \end{itemize}

    \mode<presentation>{\vfill\pause}
    \item \textmark{Operation states}
      \begin{itemize}
        \item Created when a sender is connected to a receiver, and then started.
        \item Represent the state of an asynchronous operation.
      \end{itemize}
  \end{itemize}
\end{frame}

\begin{frame}[t,fragile]{Example usage}
\begin{lstlisting}
#include <execution>
#include <iostream>

int main() {
  auto work = std::execution::just(42) | std::execution::then([](int value) {
    return value * 2;
  });

  std::execution::sync_wait(work | std::execution::then([](int result) {
    std::cout << "Result: " << result << std::endl;
  }));

  return 0;
}
\end{lstlisting}
\end{frame}

\begin{frame}[t,fragile]{Using a thread pool scheduler}
\begin{lstlisting}
#include <execution>
#include <iostream>
#include <utility>
#include "exec/static_thread_pool.hpp" // Not standard

int main() {
  auto context = stdexec::thread_pool_scheduler{}; // Not standard
  std::execution::scheduler auto sch = context.get_scheduler();

  std::execution::sender auto begin = std::execution::schedule(sch);
  std::execution::sender auto start = std::execution::then(begin, [] {
    return 42;
  });
  std::execution::sender auto multiply = std::execution::then(start, [](int value) {
    return value * 2;
  });
  std::execution::sender auto print = std::execution::then(multiply, [](int result) {
    std::cout << "Result: " << result << std::endl;
  });
  std::execution::sync_wait(std::move(print));

  return 0;
}
\end{lstlisting}
\end{frame}

\begin{frame}[t,fragile]{Simplifying notation with pipeline operator}
\begin{lstlisting}
#include <execution>
#include <iostream>
#include <utility>
#include "exec/static_thread_pool.hpp" // Not standard

int main() {
  auto context = stdexec::thread_pool_scheduler{}; // Not standard
  std::execution::scheduler auto sch = context.get_scheduler();

  std::execution::sender auto work = std::execution::schedule(sch)
    | std::execution::then([] {
        return 42;
      })
    | std::execution::then([](int value) {
        return value * 2;
      })
    | std::execution::then([](int result) {
        std::cout << "Result: " << result << std::endl;
      });
  std::execution::sync_wait(std::move(work));

  return 0;
}
\end{lstlisting}
\end{frame}